\documentclass[12pt,a4paper]{article}

\usepackage[utf8]{inputenc}
\usepackage[serbian]{babel}
\addto\captionsserbian{
    \renewcommand{\abstractname}{Сажетак}
    \renewcommand{\figurename}{Слика}
    \renewcommand{\proofname}{Доказ}
    \renewcommand{\tablename}{Табела}
}
\usepackage[T2A]{fontenc}
\usepackage{amsmath}
\usepackage{amssymb}
\usepackage{amsthm}
\usepackage{geometry}
\usepackage{xcolor}
\usepackage{graphicx}
\usepackage{wrapfig}
\usepackage{booktabs}
\usepackage{enumitem}
\usepackage{tikz}
\usetikzlibrary{positioning}

\geometry{margin=2.5cm}

\newtheorem{definicija}{Дефиниција}[section]
\newtheorem{lema}{Лема}[section]
\newtheorem{teorema}{Теорема}[section]

\title{\textbf{Основни појмови теорије графова и њихова примена}}
\author{Марко Граорац}
\date{\today}

\begin{document}



\begin{abstract}
Теорија графова представља област дискретне математике која проучава
структуре састављене од чворова и веза између њих. У овом раду представљени су основни појмови теорије графова, особине чворова и једна од основних лема позната као лема о руковању.
\end{abstract}

\section{Увод}

Графови представљају математички модел погодан за описивање односа између
различитих објеката. Објекти се представљају чворовима, док се
њихове међусобне везе представљају гранама. Због своје једноставности,
графови имају примену у рачунарству, саобраћају, телекомуникацији и
разним другим областима.

\section{Основни појмови}

\begin{definicija}
Граф је уређена двојка $G=(V,E)$, где је $V$ скуп чворова, а $E$ скуп
грана које повезују парове чворова.
\end{definicija}

Број чворова графа означава се са $|V|$, док се број његових грана
означава са $|E|$.

\begin{table}[h]
\centering
\caption{Основни појмови}
\begin{tabular}{ll}
\toprule
\textbf{Појам} & \textbf{Дефиниција} \\
\midrule
инцидентна грана & спаја два чвора \\
инцидентан чвор & спојен са одређеном граном \\
мост & његовим уклањањем раздвајају се графови \\
\bottomrule
\end{tabular}
\label{tab:vrste-grafova}
\end{table}

На следећој слици је приказан једноставан прост граф са четири чвора.

\begin{wrapfigure}{l}{0.42\textwidth}
\centering
\begin{tikzpicture}[
vertex/.style={circle, draw, minimum size=8mm}
]
\node[vertex] (A) {A};
\node[vertex, right=1.3cm of A] (B) {B};
\node[vertex, below=1.1cm of A] (C) {C};
\node[vertex, right=1.3cm of C] (D) {D};

\draw (A)--(B);
\draw (A)--(C);
\draw (B)--(D);
\draw (C)--(D);
\draw (A)--(D);
\end{tikzpicture}
\caption{Пример графа}
\label{fig:graf}
\end{wrapfigure}

Степен чвора представља број грана које су повезане са тим чвором.
На пример, на слици \ref{fig:graf}, чвор $A$ има степен $3$.

Графови се могу поделити према различитим особинама. неке од основних врста су:


\begin{itemize}
\item \textbf{прост граф:} нема петље ни вишеструке гране;
\item \textbf{усмерен граф:} гране имају одређени смер;
\item \textbf{неусмерен граф:} гране немају одређени смер;
\item \textbf{повезан граф:} посотоји пут између свака два чвора.
\end{itemize}

\section{Степени чвора}

За чвор $v\in V$ његов степен означава се са $\deg(v)$.
Степен представља број грана инцидентних са тим чвором.

На пример, за граф са слике \ref{fig:graf} важи:

$$
\deg(A)=3,\qquad
\deg(B)=2,\qquad
\deg(C)=2,\qquad
\deg(D)=3.
$$

Збир степена свих чворова је:

$$
\deg(A)+\deg(B)+\deg(C)+\deg(D)=10.
$$

Пошто свака грана доприноси по један степен сваком од своја два
крајња чвора, очекујемо да је збир степена једнак двоструком броју
грана.

\begin{lema}[Лема о руковању]
За сваки коначан неусмерен граф $G=(V,E)$ важи:

$$
\sum_{v\in V}\deg(v)=2|E|.
$$

\end{lema}

\begin{proof}
Свака грана има два крајња чвора. Због тога свака грана доприноси
укупном збиру степена тачно $2$. Ако граф има $|E|$ грана, укупан збир
степена једнак је $2|E|$.
\end{proof}

\section{Последица}

Из претходне леме непосредно следи да је број чворова непарног степена
увек паран.

\begin{teorema}
у сваком коначном неусмереном графу број чворова непарног степена је
паран.
\end{teorema}

\begin{proof}
Према леми о руковању, збир свих степена једнак је $2|E|$, па је
збир паран број. Збир парног броја непарних целих бројева је паран 
само ако је број непарних сабирака паран. Зато граф мора имати паран 
број чворова непарног степена.
\end{proof}

Овај резултат може бити користан при провери да ли одређена структура уопште
може представљати граф.


За дати граф важи:

$$
2|E|=2\cdot5=10,
$$

што је једнако збиру степена свих чворова.

Графови се у рачунарству могу користити за различите проблеме, на пример:

\begin{enumerate}[label=\arabic*.]
\item представљање рачунарских мрежа;
\item проналажење најкраћег пута;
\item моделовање друштвених мрежа;
\item представљање зависности између објеката.
\end{enumerate}

\section{Закључак}

Теорија графова пружа једноставан начин за математичко моделовање
веза између објеката. У раду су представљени основни појмови,
дефиниција грана, степен чворова, лема о руковању и теорема о парности 
броја чворова непарног степена. Иако су представљени резултати 
једноставни, они чине основу за проучавање сложенијих проблема из
теорије графова и рачунарства.

\end{document}
